% !TEX program = xelatex
% ============================================================
%  دفتر المختبر — Lab Notebook Template
%  نموذج دفتر مختبر منظّم بمداخل تجارب متعددة (Structured Experiment Entries)
%  Compile with: xelatex main.tex (run twice)
% ============================================================
%  Features:
%    - Structured experiment entries: ID, Date, Objective, Materials, Procedure
%    - Observations, Results, Notes sections per entry
%    - Multiple entry blocks using tables
%    - Fancy headers with notebook title
% ============================================================

\documentclass[11pt,a4paper]{article}

% ---- Geometry ----
\usepackage[a4paper,margin=2.5cm,top=2.5cm,bottom=2.5cm,headheight=16pt]{geometry}

% ---- Core packages ----
\usepackage{graphicx}
\usepackage{array}
\usepackage{booktabs}
\usepackage{tabularx}
\usepackage{multirow}
\usepackage{enumitem}
\usepackage{float}
\usepackage{caption}
\usepackage{titlesec}
\usepackage{fancyhdr}
\usepackage{setspace}
\usepackage{xcolor}
\usepackage{amsmath}

% ---- RTL / Arabic language support (load before hyperref) ----
\usepackage{bidi}
\usepackage[no-math]{fontspec}
\usepackage[quiet]{polyglossia}

% ---- Fonts ----
\setmainfont[Scale=1]{NotoKufiArabic-Regular.ttf}
\newfontfamily\arabicfont[Script=Arabic,Scale=0.95]{NotoKufiArabic-Regular.ttf}
\newfontfamily\englishfont[Scale=0.95]{TeX Gyre Termes}
\setmonofont{NotoKufiArabic-Regular.ttf}
\newfontfamily\arabicfontsf{NotoKufiArabic-Regular.ttf}

% ---- Languages ----
\setdefaultlanguage[calendar=gregorian,hijricorrection=1]{arabic}
\setotherlanguage[variant=british]{english}

% ---- Hyperref ----
\usepackage[breaklinks,hidelinks]{hyperref}
\hypersetup{
  pdftitle={دفتر المختبر},
  pdfauthor={اسم الباحث}
}

% ---- Captions in Arabic ----
\addto\captionsarabic{%
  \renewcommand{\tablename}{الجدول}%
  \renewcommand{\figurename}{الشكل}%
}

% ---- Section formatting ----
\titleformat{\section}{\Large\bfseries}{\thesection}{0.7em}{}
\titlespacing*{\section}{0pt}{1.4ex plus 0.4ex}{0.9ex plus 0.3ex}

% ---- Headers / footers ----
\pagestyle{fancy}
\fancyhf{}
\fancyhead[R]{\small\itshape دفتر المختبر (Lab Notebook)}
\fancyhead[L]{\small اسم الباحث}
\fancyfoot[C]{\small\thepage}
\renewcommand{\headrulewidth}{0.3pt}
\renewcommand{\footrulewidth}{0pt}

% ---- List spacing ----
\setlist[itemize]{topsep=0.3em, itemsep=0.15em, parsep=0pt}
\setlist[enumerate]{topsep=0.3em, itemsep=0.15em, parsep=0pt}

% ---- Table column types ----
\newcolumntype{L}[1]{>{\raggedright\arraybackslash}p{#1}}
\newcolumntype{C}[1]{>{\centering\arraybackslash}p{#1}}
\newcolumntype{Y}{>{\raggedright\arraybackslash}X}

% ---- Line spacing ----
\setstretch{1.4}

% ---- Experiment entry info table ----
\newcommand{\entryinfo}[2]{%
\begin{table}[H]
\centering
\begin{tabularx}{\textwidth}{L{3.5cm} X L{3.5cm} X}
\toprule
\textbf{رقم التجربة (ID):} & #1 & \textbf{التاريخ (Date):} & #2 \\
\bottomrule
\end{tabularx}
\end{table}
}

\begin{document}

% ============================================================
% COVER / HEADER
% ============================================================
\thispagestyle{fancy}

\begin{center}
{\LARGE\bfseries دفتر المختبر (Lab Notebook)}\\[0.8em]
% TODO: اكتب اسم الباحث والمختبر هنا.
{\large اسم الباحث --- المختبر، القسم، الجامعة}\\[0.3em]
{\small \today}
\end{center}

\vspace{1em}

% ============================================================
% ENTRY 1
% ============================================================
\section{التجربة الأولى}

% TODO: عدّل رقم التجربة والتاريخ.
\entryinfo{\LR{EXP-001}}{\today}

\subsection*{الهدف (Objective)}
% TODO: اكتب هدف التجربة هنا.
هذا نص تجريبي لهدف التجربة. يصف الغرض من إجراء التجربة والسؤال البحثي المطروح.

\subsection*{المواد والأدوات (Materials)}
% TODO: اذكر المواد والأدوات المستخدمة.
\begin{itemize}
    \item مادة 1 --- الكمية
    \item مادة 2 --- الكمية
    \item جهاز 1 --- المواصفات
\end{itemize}

\subsection*{طريقة العمل (Procedure)}
% TODO: اكتب خطوات العمل بالترتيب.
\begin{enumerate}
    \item الخطوة الأولى: وصف الخطوة.
    \item الخطوة الثانية: وصف الخطوة.
    \item الخطوة الثالثة: وصف الخطوة.
\end{enumerate}

\subsection*{الملاحظات (Observations)}
% TODO: اكتب الملاحظات أثناء التجربة.
هذا نص تجريبي للملاحظات. يسجل ما لوحظ أثناء إجراء التجربة.

\subsection*{النتائج (Results)}
% TODO: اكتب النتائج هنا.

\begin{table}[H]
\centering
\caption{نتائج التجربة الأولى}
\begin{tabularx}{\textwidth}{C{3cm} C{3cm} C{3cm}}
\toprule
\textbf{العينة} & \textbf{القيمة المقاسة} & \textbf{الوحدة} \\
\midrule
% TODO: املأ بيانات النتائج.
1 & \LR{10.5} & \LR{mg/L} \\
2 & \LR{12.3} & \LR{mg/L} \\
3 & \LR{9.8}  & \LR{mg/L} \\
\bottomrule
\end{tabularx}
\end{table}

\subsection*{ملاحظات إضافية (Notes)}
% TODO: اكتب أي ملاحظات إضافية أو خطوات تالية.
هذا نص تجريبي للملاحظات الإضافية. يذكر أي تحفظات أو خطوات تالية مقترحة.

\newpage

% ============================================================
% ENTRY 2
% ============================================================
\section{التجربة الثانية}

% TODO: عدّل رقم التجربة والتاريخ.
\entryinfo{\LR{EXP-002}}{\today}

\subsection*{الهدف (Objective)}
% TODO: اكتب هدف التجربة هنا.

\subsection*{المواد والأدوات (Materials)}
% TODO: اذكر المواد والأدوات المستخدمة.
\begin{itemize}
    \item مادة 1 --- الكمية
    \item جهاز 1 --- المواصفات
\end{itemize}

\subsection*{طريقة العمل (Procedure)}
% TODO: اكتب خطوات العمل بالترتيب.
\begin{enumerate}
    \item الخطوة الأولى: وصف الخطوة.
    \item الخطوة الثانية: وصف الخطوة.
\end{enumerate}

\subsection*{الملاحظات (Observations)}
% TODO: اكتب الملاحظات أثناء التجربة.

\subsection*{النتائج (Results)}
% TODO: اكتب النتائج هنا.

\begin{table}[H]
\centering
\caption{نتائج التجربة الثانية}
\begin{tabularx}{\textwidth}{C{3cm} C{3cm} C{3cm}}
\toprule
\textbf{العينة} & \textbf{القيمة المقاسة} & \textbf{الوحدة} \\
\midrule
% TODO: املأ بيانات النتائج.
1 & \LR{22.1} & \LR{mg/L} \\
2 & \LR{24.0} & \LR{mg/L} \\
\bottomrule
\end{tabularx}
\end{table}

\subsection*{ملاحظات إضافية (Notes)}
% TODO: اكتب أي ملاحظات إضافية.

\newpage

% ============================================================
% ENTRY 3 (template — copy as needed)
% ============================================================
\section{التجربة الثالثة}

% TODO: انسخ كتلة التجربة هذه حسب الحاجة لإضافة مداخل جديدة.
\entryinfo{\LR{EXP-003}}{\today}

\subsection*{الهدف (Objective)}
% TODO: اكتب هدف التجربة هنا.

\subsection*{المواد والأدوات (Materials)}
% TODO: اذكر المواد والأدوات المستخدمة.

\subsection*{طريقة العمل (Procedure)}
% TODO: اكتب خطوات العمل بالترتيب.

\subsection*{الملاحظات (Observations)}
% TODO: اكتب الملاحظات أثناء التجربة.

\subsection*{النتائج (Results)}
% TODO: اكتب النتائج هنا.

\subsection*{ملاحظات إضافية (Notes)}
% TODO: اكتب أي ملاحظات إضافية.

\end{document}
